% =============================================================================
% specforge -- TikZ style library for block diagrams and timing diagrams.
% colour-guide.md Sections 5.4, 5.7, 12 and 13. Colours and line widths come
% from build/sf-config.tex (styles/style-defaults.md, groups diagrams and
% timing); shapes and arrow tips are fixed here.
%
% Block diagrams
%   boundary      top-level module boundary            block     functional block
%   ifgroup       interface-group box                  port      port square on a boundary
%   wire          plain connection                     sigarrow  connection with arrow
%   ctrl/ctrlarrow  control signal (thin)              data/dataarrow  data or state (thick)
%   bundle/bundlearrow  signal bundle or information flow  hlarrow  highlighted path (label it!)
%   packet/packetarrow  one packet (struct, interface): one arrow, no fields
%   junction, djunction  wire junctions (thin, thick)  zoomline  zoom guide (dashed)
%   siglabel  signal name     flowlabel  information in words
%   blockname / modname  block / module name
%   blocksub  block sub-label grouptitle           group title     widthlabel  bus width
% State machines (see latex-document rule 04)
%   fsmstate   state: a circle, never a box       fsmresetdot, fsmreset  reset entry
%   fsmtrans   state transition (straight, bend left/right, loop above/...)
%   fsmcond    transition label: condition / outputs
%   fsminput, fsmoutput  signals entering or leaving the machine (dashed)
%   \fsmout{text}  outputs inside a state circle (below the name)
%   \fsmlegend{(x,y)}, \fsmlegendio{(x,y)}  legend of the arrow kinds
% Buffer and pointer illustrations
%   entry, entryhl (highlighted), entrydim (invalid), entrywarn (warning; label it)
%   ptrarrow  pointer arrow   rangebrace  brace under a range of entries
% Timing diagrams (tikz-timing)
%   style sfwave on the tikztimingtable; \wavecycles{n} draws the cycle grid
%   and numbers, \wavenote{first}{last}{text} adds a note under cycles.
%   One clock cycle is two timing units wide; the clock row is [sfclock]4{HL}.
%   Ragged-right text distorts tikz-timing tables: the build wraps every
%   figure named fig-wave-* or fig-timing-* in justified text.
% =============================================================================
\usetikzlibrary{arrows.meta,calc,positioning,patterns,backgrounds,
                decorations.pathreplacing}

% ---- Timing diagrams ----------------------------------------------------------
% Font of size #1 pt with 1.2 x leading.
\newcommand{\sfsize}[1]{\fontsize{#1}{\fpeval{1.2*(#1)}}\selectfont}
\newcommand{\wavefont}{\ttfamily\sfsize{\sfTimingLabelSizePt}}
\tikzset{
  sfwave/.style={
    xscale=\sfTimingXScale, yscale=\sfTimingYScale,
    line width=\sfTimingLinePt pt, color=\sfTimingSignalColour,
    timing/slope=\sfTimingSlope,
    timing/rowdist=2,
    timing/name/.style={font=\wavefont, text=\sfTimingSignalColour},
    timing/d/text/.style={font=\wavefont, text=\sfTimingSignalColour},
    timing/d/background/.style={fill=bgpage},
    timing/u/.style={draw=\sfTimingClockColour},
    timing/u/background/.style={pattern=north east lines,
                                pattern color=\sfTimingUndefinedPatternColour},
  },
  % Row option for the clock: [sfclock]4{HL}
  sfclock/.style={color=\sfTimingClockColour},
  wavecycle/.style={font=\sfsize{\sfTimingLabelSizePt},
                    text=textsecondary, inner sep=1pt},
  wavenote/.style={font=\sfsize{\sfTimingLabelSizePt},
                   text=textprimary, inner sep=1pt, align=center},
}

% Cycle grid and cycle numbers. #1 = number of cycles shown.
\newcommand{\wavecycles}[1]{%
  \begin{scope}[on background layer]
    \foreach \c in {0,...,#1} {
      \draw[\sfTimingGridColour, line width=\sfTimingGridLinePt pt]
        (2*\c,1.45) -- (2*\c,-2*\nrows+1.75);
    }
  \end{scope}
  \foreach \c in {1,...,#1} {
    \node[wavecycle] at (2*\c-1,1.85) {\the\numexpr\c-1\relax};
  }
  \node[wavecycle, anchor=east] at (-0.35,1.85) {cycle};
}

% Text note under the waveform. #1 = first cycle, #2 = last cycle, #3 = text.
\newcommand{\wavenote}[3]{%
  \node[wavenote, anchor=north] at (#1+#2+1,-2*\nrows+1.6) {#3};
}

% ---- Block diagrams -----------------------------------------------------------
\tikzset{
  boundary/.style={draw=\sfDiagramsBoundaryDraw, fill=\sfDiagramsBoundaryFill,
                   line width=\sfDiagramsBoundaryLinePt pt},
  block/.style={draw=\sfDiagramsBlockDraw, fill=\sfDiagramsBlockFill,
                line width=\sfDiagramsBlockLinePt pt, rectangle},
  ifgroup/.style={draw=\sfDiagramsGroupDraw, fill=\sfDiagramsGroupFill,
                  line width=\sfDiagramsGroupLinePt pt,
                  rounded corners=\sfDiagramsGroupCornerRadiusPt pt},
  wire/.style={draw=\sfDiagramsWireColour, line width=\sfDiagramsWireLinePt pt,
               line cap=round, line join=round},
  sigarrow/.style={wire, -{Stealth[length=2.3mm,width=1.7mm]}},
  ctrl/.style={draw=\sfDiagramsWireColour, line width=\sfDiagramsControlWirePt pt,
               line cap=round, line join=round},
  ctrlarrow/.style={ctrl, -{Stealth[length=2.2mm,width=1.6mm]}},
  data/.style={draw=\sfDiagramsWireColour, line width=\sfDiagramsDataWirePt pt,
               line join=round},
  dataarrow/.style={data, -{Stealth[length=3.2mm,width=2.6mm]}},
  bundle/.style={draw=\sfDiagramsWireColour, line width=\sfDiagramsBundlePt pt,
                 line cap=butt},
  bundlearrow/.style={bundle, -{Stealth[length=2.8mm,width=2.2mm]}},
  % one packet (struct, interface, field group) as a whole: double line,
  % open head; its fields are never drawn
  packet/.style={draw=\sfDiagramsWireColour, line width=\sfDiagramsPacketLinePt pt,
                 double=bgpage, double distance=\sfDiagramsPacketGapPt pt, line cap=butt},
  packetarrow/.style={packet, -{Latex[open, length=3.4mm, width=3.2mm, line width=\sfDiagramsPacketLinePt pt]}},
  hlarrow/.style={draw=\sfDiagramsHighlightColour,
                  line width=\sfDiagramsHighlightLinePt pt, line join=round,
                  -{Stealth[length=2.5mm,width=1.9mm]}},
  port/.style={draw=\sfDiagramsBoundaryDraw, fill=bgpage,
               line width=\sfDiagramsWireLinePt pt, rectangle,
               minimum size=\sfDiagramsPortSizeMm mm, inner sep=0pt},
  junction/.style={circle, fill=\sfDiagramsWireColour, draw=none,
                   minimum size=\sfDiagramsJunctionSizeMm mm, inner sep=0pt},
  djunction/.style={junction, minimum size=1.5*\sfDiagramsJunctionSizeMm mm},
  zoomline/.style={draw=\sfDiagramsBoundaryDraw, line width=0.7pt,
                   dash pattern=on 3pt off 2pt},
  siglabel/.style={font=\ttfamily\sfsize{\sfDiagramsSignalLabelSizePt},
                   text=textprimary, inner sep=1.6pt},
  blockname/.style={font=\headingfont\sfsize{\sfDiagramsBlockNameSizePt},
                    text=textprimary, align=center},
  modname/.style={font=\headingfont\sfsize{\fpeval{\sfDiagramsBlockNameSizePt-1}},
                  text=textprimary, align=center},
  blocksub/.style={font=\sfsize{\sfDiagramsBlockSubSizePt},
                   text=textsecondary, align=center},
  grouptitle/.style={font=\bfseries\sfsize{\sfDiagramsGroupTitleSizePt},
                     text=textprimary, inner sep=1.6pt},
  % what an information flow carries, in words (not a signal name)
  flowlabel/.style={font=\sfsize{\sfDiagramsSignalLabelSizePt},
                    text=textprimary, inner sep=1.6pt},
  widthlabel/.style={font=\sfsize{\sfDiagramsWidthLabelSizePt},
                     text=textsecondary, inner sep=1.6pt},
}

% ---- Buffer and pointer illustrations ----------------------------------------
\tikzset{
  entry/.style={draw=\sfDiagramsBlockDraw, fill=bgpage,
                line width=\sfDiagramsWireLinePt pt, rectangle,
                minimum width=9mm, minimum height=7mm, inner sep=0pt,
                font=\wavefont, text=textprimary},
  entryhl/.style={entry, draw=\sfDiagramsHighlightColour, fill=accentsubtle,
                  line width=1.2pt},
  entrydim/.style={entry, fill=bgsubtle, draw=bordersubtle, text=textsecondary},
  entrywarn/.style={entry, draw=warnfg, fill=warnbg, line width=1.2pt},
  ptrarrow/.style={wire, -{Stealth[length=2.3mm,width=1.7mm]}},
  rangebrace/.style={draw=\sfDiagramsWireColour, line width=\sfDiagramsWireLinePt pt,
                     decorate, decoration={brace, amplitude=5pt, mirror}},
}

% ---- State machines -------------------------------------------------------------
% States are circles. Transitions are filled-head arrows, straight or bent;
% the reset entry starts at a filled dot; signals into or out of the
% machine are dashed. Colours and sizes come from the state_machines keys.
\tikzset{
  fsmstate/.style={circle, draw=\sfStateMachinesStateDraw, fill=\sfStateMachinesStateFill,
                   line width=\sfStateMachinesStateLinePt pt,
                   minimum size=\sfStateMachinesStateSizeMm mm, inner sep=1pt, align=center,
                   font=\ttfamily\sfsize{\sfStateMachinesNameSizePt}, text=textprimary},
  fsmtrans/.style={draw=\sfStateMachinesTransitionColour,
                   line width=\sfStateMachinesTransitionLinePt pt,
                   -{Stealth[length=2.4mm,width=1.8mm]}, shorten >=0.4pt,
                   every loop/.style={looseness=6, min distance=8mm}},
  fsmresetdot/.style={circle, fill=\sfStateMachinesResetColour, draw=none,
                      minimum size=2.2mm, inner sep=0pt},
  fsmreset/.style={draw=\sfStateMachinesResetColour, line width=\sfStateMachinesResetLinePt pt,
                   -{Stealth[length=2.8mm,width=2.2mm]}, shorten >=0.4pt},
  fsmcond/.style={font=\ttfamily\sfsize{\sfStateMachinesLabelSizePt}, text=textprimary,
                  fill=bgpage, inner sep=1pt, align=center},
  fsminput/.style={draw=\sfDiagramsWireColour, line width=\sfDiagramsWireLinePt pt,
                   dash pattern=on 3pt off 2pt, -{Stealth[length=2.2mm,width=1.6mm]}},
  fsmoutput/.style={fsminput},
}
% Outputs inside a state circle, below the name (Moore outputs).
\newcommand{\fsmout}[1]{\\{\color{textsecondary}\sfsize{\sfStateMachinesLabelSizePt}#1}}
% Legend: reset entry and transition; #1 = top left corner, e.g. (0,-3).
\newcommand{\fsmlegend}[1]{%
  \begin{scope}[shift={#1}]
    \node[fsmresetdot] (fsmlgdot) at (0,0) {};
    \draw[fsmreset] (fsmlgdot) -- (0.9,0);
    \node[flowlabel, anchor=west] at (1.0,0) {reset};
    \draw[fsmtrans] (0,-0.55) -- (0.9,-0.55);
    \node[flowlabel, anchor=west] at (1.0,-0.55) {transition: condition / outputs};
  \end{scope}}
% Legend with the dashed inputs and outputs of the machine as well.
\newcommand{\fsmlegendio}[1]{%
  \fsmlegend{#1}%
  \begin{scope}[shift={#1}]
    \draw[fsminput] (0,-1.1) -- (0.9,-1.1);
    \node[flowlabel, anchor=west] at (1.0,-1.1) {input or output signal};
  \end{scope}}
